% load package if document contains at least one citation
\usepackage[maxbibnames=99,sorting=none,style=numeric,backend=biber]{biblatex}
% enable option to print bibliography as a numbered subsection
\defbibheading{bibsubsection}[\bibname]{%
  \subsection{#1}%
}
\addbibresource{Private/PHBern/References.bib}

\begin{document}

\lhead[ \leftmark ]{PHBern, LUBM}
\rhead[Informatik]{Cyril Wendl, \today, Bern}

\part*{\acrshort{LNW} Teil 1: Schriftliche Reflexion\\\enquote{Berufsbildung im digitalen Zeitalter}}

\section{Einführung}
Während des  Seminars \enquote{\gls{LUBM}}, welches an der PHBern im Rahmen des Lehrdiplom-Studiengangs am \gls{IS2} angeboten wird, wurden verschiedene Thematiken im Zusammenhang mit dem Lernen und Unterrichten an \gls{BMS} vertieft. Im Rahmen des Seminars wurde insbesondere auf neuere Herausforderungen im Zusammenhang mit digitalen Methoden an \gls{BMS} eingegangen. Dabei ging es insbesondere um folgende Fragen:

\begin{itemize}
    \item Wie können angehende \glspl{LP} aktuelle Trends in der Berufswelt erkennen, insbesondere hinsichtlich \gls{KI}?
    \item Wie können Folgen der Digitalisierung für die Berufsbildung abgeschätzt werden?
    \item Wie können \glspl{LP} unterschiedliche Tools, Methoden und didaktische Settings sinnvoll einsetzen?
    \item Welche Richtlinien sollten betreffend der Verwendung von \gls{KI} in schulischen Arbeiten gesetzt werden?
    \item Wie verändert \gls{KI} die Rolle als \gls{LP}? Braucht es diese überhaupt noch?
\end{itemize}

Da ich bereits seit eineinhalb Jahren an der Kantonsschule im Lee unterrichte und auch hinsichtlich meines potentiellen \gls{BM}-Fachs \enquote{\gls{TU}}, welches ich mit dem Zweifächer-Lehrdiplom in Geografie und Informatik mit Zusatzqualifikation für Berufsmaturitätsschulen unterrichten könnte, interessiere ich mich stark für den sinn- und verantwortungsvollen Einsatz von digitalen Lern-Methoden im Unterricht. Aus diesen Gründen werde ich in dieser Arbeit insbesondere den verantwortungsvollen Einsatz von digitalen Lern-Methoden, insbesondere \gls{KI}, untersuchen.

Um das Thema inhaltlich einzugrenzen und zu strukturieren, werde ich mich in dieser Arbeit auf zwei konkrete Fragestellungen im Zusammenhang mit \gls{KI} konzentrieren:

\begin{itemize}
    \item Welche Voraussetzungen müssen gegeben sein, damit kompetenzorientiertes Lernen möglich ist, und wie können diese Voraussetzungen durch \gls{KI} unterstützt werden?
    \item Wie verändern \gls{KI} und insbesondere \glspl{LLM} wie ChatGPT meine Rolle als \gls{LP} und wie kann ich \gls{KI} einsetzen, um kompetenzorientierten Unterricht zu fördern?
\end{itemize}

\section{Voraussetzungen für kompetenzorientiertes Lernen und Unterstützung von Lernprozessen durch \gls{KI}}

Im Fach \gls{TU} werden sowohl quantitativ-naturwissenschaftliche Inhalte vermittelt (z.B. \enquote{Materialflüsse}, \enquote{Energie- und Stoffflüsse}) wie auch Inhalte eher qualitativ-geisteswissenschaftlicher Art (z.B. \enquote{Die Welt: ein vernetztes System} oder \enquote{Lösungsansätze zu einer nachhaltigen Entwicklung}) \cite{sbfi}. Es scheint daher naheliegend, dass die Breite dieses Fachs sich ebenfalls in den Lern- und Prüfungsaufgaben widerspiegelt, welche je nach Thema Aufgaben beinhalten, welche mehr in Richtung \enquote{Berechnung von Zahlen} gehen (mit einer klaren wahr/falsch-Antwort) oder mehr in Richtung qualitativer Antworten wie beispielsweise dem Erstellen eines Wirkungsgefüges oder einer kurzen Text-Antwort zu einer gegebenen Problematik. In beiden Fällen wird durch den Rahmenlehrplan festgehalten, dass im Zentrum der Berufsausbildung die Entwicklung fachlicher sowie überfachlicher \textit{Kompetenzen} steht \cite[7-8]{sbfi}. \enquote{Träges} Wissen soll also vermieden werden und der kompetenzorientierte Unterricht soll authentischer und problembasierter gestaltet werden, damit sich die Lernenden den komplexen und rasch wandelnden Umgebungen ihres Betriebs anpassen können, wie dies etwa auch von \textcite[107-117]{alma99117272661305511} gefordert wird.

Beide Typen von Lernaufgaben können durch aktuelle \gls{KI}-Sprachmodelle wie beispielsweise ChatGPT relativ erfolgreich erledigt werden, wobei das Berechnen von Zahlen je nach Fall ein spezielles Prompt-Design erforderlich macht, welches jedoch mit wenig Aufwand erlernt werden kann \cite{bfh}. Aufgrund der drohenden Gefahr, der Sinn von Lernaufgaben sowie Prüfungen könnte durch die Verwendung von \gls{KI}-Tools wie ChatGPT untergraben werden, haben verschiedene Institutionen Richtlinien zur Verwendung von \gls{KI} im akademischen Kontext verfasst, wie beispielsweise die \textcite{bfh} und das \textcite{bbz}. Die Verunsicherung vieler \glspl{LP} angesichts der neuartigen digitalen Werkzeuge reflektiert sich insbesondere in Ängsten hinsichtlich Prüfungen, wie dargelegt beispielsweise in \cite{bund}. Weniger Aufmerksamkeit wurde bisher der Frage gewidmet, inwiefern \gls{KI} den kompetenzorientierten Unterricht fördert bzw. gefährdet. Um diese Frage zu beantworten, sollten jedoch zuerst die Grundvoraussetzungen für kompetenzorientiertes Lernen beschrieben werden.

Im Artikel \enquote{Lernen ist (mehr als) Arbeiten} stellt sich \textcite{nzzfolio} auf den Standpunkt, dass Lernen nur möglich ist durch Arbeit in verschiedenen Regionen des Hirns. \textcite{nzzfolio} verweist dabei auf das Kompetenz-Ressourcen-Modell, welches im Bericht von \textcite{sbfi2} verwendet wird. Laut diesem Modell wird für den Aufbau von Handlungsfähigkeit die Entwicklung der Bereiche Kenntnisse, Fähigkeiten / Fertigkeiten sowie Haltungen vorausgesetzt. \textcite{nzzfolio} verbindet jeden dieser Bereiche mit spezifischen Hirnarealen und verweist gleichzeitig auf Erfolgsfaktoren, um Lernen in den jeweiligen Hirnarealen zu ermöglichen:
\begin{itemize}
    \item Neue \textbf{Kenntnisse} müssen vom Kurzzeitgedächtnis (Hippocampus) aufgenommen und ins Langzeitgedächtnis (Hirnrinde) verschoben werden können. Dieser Vorgang wird insbesondere durch Ruhezeiten begünstigt und verstärkt, insbesondere durch Schlaf. Diese Erkenntnis wurde in letzter Zeit auch in zahlreichen Sachbüchern festgehalten, wie etwa in \parencite{alma99116800416405511}.
    \item \textbf{Fertigkeiten} werden durch viele Wiederholungen und Bewegungs-Abläufe automatisiert, wofür insbesondere das Kleinhirn zuständig ist. Auch diese Erkenntnis findet sich in zahlreichen öffentlichen Publikationen wieder \cite{huberman}.
    \item \textbf{Fähigkeiten} betreffen \enquote{das Verstehen von Zusammenhängen und das Anwenden von Methoden} \parencite{nzzfolio}, was primär durch das Erledigen von variierenden Aufgaben erreicht wird. Hier sind insbesondere manuelle Aufgaben eine wichtige Möglichkeit, durch Abwechslung weitere Areale des Gehirns zu stimulieren.
    \item Das Ausbilden von \textbf{Haltungen} wird insbesondere durch den sozialen Bereich des Gehirns (Frontallappen) ermöglicht, welcher sich erst mit der Pubertät vollständig entwickelt.
\end{itemize}

Die Kernaussage des Artikels besteht darin, dass Lernen zwar durch didaktisch sinnvoll aufbereitete digitale Lernmittel unterstützt wird, effektives Lernen jedoch erst durch eigenständige, \enquote{harte} Arbeit ermöglicht wird -- schliesslich benötigt das Gehirn im Schnitt ca. 20\% der gesamten Körperenergie. Nicht zuletzt sei Lernen massgeblich beeinflusst durch positive Emotionen, welche, bedingt durch die Evolution des Menschen als soziales Tier, primär über menschliche Kontakte stattfinden.

In Bezug auf die Anwendung von \gls{KI} in kompetenzorientiertem Unterricht kann also grundsätzlich Folgendes festgehalten werden:

\begin{itemize}
    \item Lernen erfolgt in verschiedenen Hirnarealen und kann grundsätlich nicht an ein Tool \enquote{externalisiert} werden -- Lernen erfordert also immer \enquote{harte} Arbeit durch die Lernenden selbst.
    \item Lernen kann kaum ohne emotionale Verbindungen mit dem Erlernten erfolgen. Hier sind die Lehrpersonen und Mitlernenden zentral, da sie positive (im wünschenswerteren Fall) oder negative Emotionen zum Inhalt vermitteln können. Dies ist insbesondere dann wichtig, wenn es um das Erstellen positiver Bezüge zu den zu erlernenden Inhalten geht.
    \item \gls{KI}-Tools können grundsätzlich hilfreich sein, um abwechslungsreiche Inhalte zu erstellen oder das Gelernte auf neue Weise zu vermitteln oder zu prüfen, wodurch insbesondere Fähigkeiten und Fertigkeiten gestärkt werden können. Dies ist insbesondere dann sinnvoll, wenn \gls{KI}-gestützte Inhalte auf didaktisch sinnvolle Weise (z.B. Reihenfolge, Aufbau der Aufgaben etc.) im Unterricht eingesetzt werden. Diese Erkenntnis wird ebenfalls von verschiedenen Akademischen Institutionen festgehalten, welche alle den Standpunkt vertreten, dass der Einsatz von \gls{KI} grundsätzlich gewinnbringend sein kann und daher nicht im Versteckten erfolgen, sondern aktiv thematisiert und integriert werden soll \cite{bbz,bfh}.
\end{itemize}

Nebst den genannten Möglichkeiten, welche durch den Einsatz von \gls{KI} im Unterricht zur Verfügung gestellt werden, gilt es, sich deren Limitationen vor Auge zu halten, da gerade hinsichtlich didaktischer Aspekte zum aktuellen Zeitpunkt noch viele Fragen offen sind. Die \textcite{bfh} weist in diesem Zusammenhang unter anderem auf folgende derzeitige Limitationen von \glspl{LLM} wie ChatGPT hin:

\begin{itemize}
    \item Der Text, der generiert wird, soll möglich plausibel klingen, allerdings ist ChatGPT nicht imstande, die faktuell-inhaltliche Richtigkeit des generierten Texts zu überprüfen.
    \item ChatGPT kann nicht garantieren, dass ein generierter Text einen Bezug zur echten Welt hat: Beispielsweise können Referenzen in der Bibliographie erfunden sein.
    \item ChatGPT kann Lösungswege für Probleme beschreiben, häufig kann es jedoch den Lösungsweg für eigene Berechnungen und Lösungen nicht angeben.
    \item ChatGPT funktioniert, wie viele andere Algorithmen auch, nach dem \gls{GIGO}-Prinzip, welches besagt, dass die Qualität der Antworten von ChatGPT nur so gut ist wie die Qualität der Trainingsdaten. ChatGPT kann also möglicherweise auch Vorurteile perpetuieren durch \textit{biases} in den Trainingsdaten.
    \item Gerade der menschlich anmutende Aspekt moderner \gls{KI}-Tools wie ChatGPT, welche ihre eigenen Aussagen relativieren und mit echten Menschen \enquote{Konversationen} führen können, verleitet laut der \textcite{bfh} dazu, Tools wie ChatGPT allzu viel Menschlichkeit zuzuschreiben, was im schlimmsten Falle zu blindem Vertrauen führen kann.
\end{itemize}

Weitere Limitierungen von ChatGPT sind die häufig oberflächlichen Antworten, mangelnde Anpassung an persönliche Präferenzen (Stil, spezielle Wörter in gewissen Sprachräumen wie etwa der Schweiz, etc.) und mangelnde Kenntnisse des Kontexts. Obwohl viele dieser Limitationen durch geschicktes \textit{prompt engineering}, also die Verfeinerung der Abfrage, überwunden werden können, erfordert dies oft tiefere Kenntnisse der Tools und häufige Anpassungen an sich ständig ändernde Benutzeroberflächen und \gls{KI}-Versionen.

Eine Hauptlimitation von ChatGPT scheint zudem die mangelnde Reproduzierbarkeit von Antworten: Dies bedeutet, dass Antworten immer anders ausfallen, jedes Mal wenn die Abfrage abgeschickt wird, selbst wenn der Abfragetext immer in derselben Form verfasst wird. Dies erschwert den Einsatz von ChatGPT im Unterricht, wo eine gewisse Standardisierung der Inhalte wünschenswert ist, um eine allzu grosse Heterogenität der Wissens- und Fähigkeitsniveaus innerhalb der Klasse zu vermeiden.

\section{Meine Rolle als \gls{LP} und Einsatz von \gls{KI} im eigenen Unterricht}
Kompetenzorientiertes Lernen erfordert harte Arbeit durch die Lernenden und die Aktivierung unterschiedlicher Areale. Diese kann grundsätzlich durch \gls{KI}-Tools nicht abgenommen werden, allerdings können \gls{KI}-Tools unterschiedliche Möglichkeiten bereitstellen, wie etwa die Vereinfachung oder Umformulierung von Texten, die Erstellung von Aufgabenstellungen, Lückentexten oder Glossaren, die Unterstützung bei der Erstellung von Leistungskontrollen, automatisierte sprachliche Korrekturen, Impulse bei der Planung von Unterrichtssequenzen sowie die Entwicklung neuer, kreativer Inhalte \cite{bbz}. Die Erstellung von Lückentexten könnte etwa in qualitativ orientierten Teilbereichen des Fachs \gls{TU} wie dem Bereich \enquote{1.2 -- Vernetzte Systeme} lernförderlich wirken, währenddem \gls{KI}-Tools gerade bei der kreativen Bereicherung von Unterrichtssequenzen zu eher theoretischen Themen wie beispielsweise \enquote{2.2 -- Energie und Energieflüsse} nützlich sein könnten.

Limitationen aktueller \gls{KI}-Tools werfen die Frage auf, inwiefern volatile, teilweise falsche und nicht reproduzierbare Ergebnisse von \gls{KI}-Tools die Qualität des Unterrichts potentiell verschlechtern können. Eigene Erfahrungen bei der Erstellung von Lernaufgaben im Programmierunterricht haben mir beispielsweise aufgezeigt, dass die Lösungsvorschläge häufig syntaktische Elemente enthalten, die den Lernenden noch nicht geläufig sind. Aus all diesen Gründen wird breit empfohlen, dass \glspl{LP} sowie Studierende die Inhalte immer zuerst kritisch hinterfragen, bevor sie diese in eigenen Arbeiten oder im Unterricht einsetzen \cite{bbz,bfh}. Nicht zuletzt sollte der geschickte Umgang mit \gls{KI} durch \glspl{LP} sowie durch Auszubildende erlernt werden, nicht nur um die Limitationen von \gls{KI}-Tools besser zu verstehen, sondern auch, um diese gewinnbringend einzusetzen \cite{nzzAS,nzzChatGPT}.

\section{Konklusion}
Nebst dem Potential von \gls{KI}, den Unterricht variierter, klarer und eventuell auch kreativer zu gestalten, gilt es festzuhalten, dass Lernen immer noch ein Prozess ist, welcher von Mensch zu Mensch abläuft. Emotionen spielen eine zentrale Rolle im Lernprozess und können zum aktuellen Zeitpunkt nur bedingt durch Maschinen vermittelt werden \cite{nzzfolio}. Dass dies auch in naher Zukunft so bleiben wird, scheint durch die evolutive Prägung des Menschen als soziales, in Gesellschaft agierendes Wesen naheliegend. Aus diesen Gründen sollten bei aller Euphorie für die neuen Möglichkeiten von \gls{KI} insbesondere folgende Punkte nicht vergessen gehen: \begin{enumerate*}\item Die oberflächlichen und teilweise verkürzten Antworten von \glspl{LLM} könnten insbesondere hinsichtlich kompetenzorientiertem Unterricht Schaden anrichten, gerade wenn die Antworten aufgrund ihres \enquote{menschlichen} Stils allzu ernst genommen werden. \item Die Vermittlung kompetenzorientierten Unterrichts bedarf stets einer Person, welche die \gls{KI}-Tools kritisch hinterfragt und mit ihrem eigenen Expertenwissen anreichert.
\end{enumerate*} Gerade hinsichtlich dem letzten Punkt scheint es, dass \glspl{LP} einen Beitrag zur Reifung der Lernenden leisten können, sofern sie nebst ihrem Fachwissen auch das Wissen zum verantwortungs- und sinnvollen Umgang mit \gls{KI}-Tools vermitteln. Wenn all diese Aspekte gemeinsam berücksichtigt werden, kann \gls{KI} durchaus ihren Beitrag leisten zu einem effektiveren kompetenzorientierten Unterricht.


\printglossary

\printbibliography

\end{document}